\section{Trabajos Relacionados}
\label{sec:trabajos-relacionados}

\subsection{Nota de honestidad metodológica sobre el rigor de esta revisión}

La guía exige una búsqueda sistemática siguiendo la disciplina PRISMA 2020 \cite{page2021}. Se declara
explícitamente, antes de presentar los resultados, \textbf{qué tan rigurosa fue realmente esta búsqueda}: no
se tuvo acceso directo a bases de datos académicas indexadas (Scopus, IEEE Xplore, ACM Digital Library)
con capacidad de exportar conteos exactos de resultados por cadena de búsqueda; la búsqueda se realizó
mediante un motor de búsqueda web general, en 24 consultas dirigidas durante la elaboración de este
informe, complementadas con la técnica de \textit{snowballing} de Wohlin \cite{wohlin2014ease} (seguir
referencias desde los trabajos encontrados). Esto \textbf{no es una revisión sistemática de literatura
completa} en el sentido estricto de Kitchenham y Charters \cite{kitchenham2007} --- no hay protocolo
pre-registrado, no hay doble revisor independiente, y la cobertura de bases de datos es incompleta. Se
presenta como lo que realmente es: una búsqueda dirigida y honesta, reportada con la disciplina de
PRISMA en la medida de lo posible, no como una SLR formal. Esta limitación se declara también en la
Sección~\ref{sec:amenazas-validez}. Las palabras clave usadas, sin cadena booleana exacta ni fecha
verificable por consulta, se listan en el Anexo A (Sección~\ref{anexo:busqueda}).

\subsection{Diagrama de flujo (adaptado de PRISMA 2020)}

La Figura~\ref{fig:prisma} resume las cuatro etapas de esa búsqueda dirigida: identificación, cribado,
elegibilidad e inclusión en la comparación directa.

\begin{figure}[H]
\centering
\includegraphics[width=0.75\textwidth]{figuras/fig-prisma-flujo.pdf}
\caption{Literature search flow diagram, adapted from PRISMA 2020 \cite{page2021} to this team's real
database-access limitations. Fuente Mermaid versionada en \texttt{docs/checklists/PRISMA-2020.md}
(mismo mecanismo de renderizado que los 3 diagramas C4 del Capítulo 6), a partir de las mismas cifras
citadas en el cuerpo del texto — 24 consultas, 32 trabajos identificados, 32 elegibles, 13 incluidos en
la comparación directa (Tabla~\ref{tab:trabajos-relacionados}).}
\label{fig:prisma}
\end{figure}

\subsection{Por qué no hay trabajos previos sobre ``sistemas de gestión de pre-sustentaciones''}

La búsqueda dirigida no encontró literatura académica indexada específicamente sobre sistemas de software
para gestionar el proceso de pre-sustentación o defensa de tesis de pregrado como fenómeno de
investigación en sí mismo --- es un dominio de aplicación institucional específico (UTEQ, y de forma más
amplia, universidades ecuatorianas), no un tema de investigación en ingeniería de software per se. En su
lugar, la comparación de la Tabla~\ref{tab:trabajos-relacionados} se centra en los trabajos empíricos que
fundamentan las \textbf{técnicas y métodos} usados en este proyecto (autenticación JWT, integración
continua, análisis de dependencias vulnerables, revisión de código, trazabilidad de requisitos) ---
que es la comparación honesta y relevante que la evidencia disponible permite hacer.

\subsection{Tabla comparativa}

\begin{table}[H]
\centering
\small
\begin{tabular}{p{2.6cm}p{2cm}p{1.8cm}p{6.2cm}}
\toprule
\textbf{Trabajo} & \textbf{Dominio} & \textbf{Venue} & \textbf{Relación con este proyecto} \\
\midrule
Vasilescu et al. \cite{vasilescu2015} & CI/GitHub & FSE 2015 & Fundamenta la decisión de CI automatizado (\texttt{ci.yml}, Sección~\ref{sec:implementacion}) \\
Hilton et al. \cite{hilton2016} & CI open-source & ASE 2016 & Costos/beneficios de CI cuantificados; contrasta con la ausencia de CI hasta la Fase 5 de este proyecto \\
Zampetti et al. \cite{zampetti2020} & Malas prácticas de CI & Empirical SE 2020 & Catálogo de 79 anti-patrones usado para revisar \texttt{ci.yml} \\
Beller et al. \cite{beller2017} & Fallos de build en CI & MSR 2017 & Motiva el diseño no-bloqueante del paso de análisis estático en \texttt{ci.yml} \\
Pashchenko et al. \cite{pashchenko2018} & Dependencias vulnerables & ESEM 2018 & Marco para interpretar los hallazgos de \texttt{npm audit} (Sección~\ref{sec:evaluacion}) \\
Decan et al. \cite{decan2018} & Vulnerabilidades npm & MSR 2018 & Contexto cuantitativo del ecosistema npm citado en la discusión de dependencias \\
Rezaei Nasab et al. \cite{rezaeinasab2022} & Seguridad en microservicios & J. Systems and Software 2022 & 28 prácticas de seguridad contrastadas contra las implementadas (JWT, rate limiting, cabeceras) \\
Lu et al. \cite{lu2018} & Rate limiting de autenticación & ACSAC 2018 & Halló que 72\,\% de sitios reales no limitan intentos de login; motiva la implementación de rate limiting de este proyecto \\
Bacchelli \& Bird \cite{bacchelli2013} & Revisión de código & ICSE 2013 & Contrasta con la ausencia de un proceso formal de code review en este proyecto (Sección~\ref{sec:amenazas-validez}) \\
Mäder \& Egyed \cite{mader2015} & Trazabilidad de requisitos & Empirical SE 2015 & Fundamenta el valor de \texttt{matriz.csv} más allá de un requisito administrativo \\
McIntosh et al. \cite{mcintosh2016} & Revisión de código y calidad & Empirical SE 2016 & Contexto para la ausencia de un proceso formal de code review en este proyecto (Sección~\ref{sec:amenazas-validez}) \\
Maximilien \& Williams \cite{maximilien2003} & TDD y reducción de defectos & ICSE 2003 & Contrasta con la cobertura de 82.01\,\% de líneas de este proyecto (Sección~\ref{sec:evaluacion}); TDD no se practicó formalmente \\
Gallaba et al. \cite{gallaba2022} & Datos operacionales de CI & ICSE 2022 & Fundamenta el diseño del pipeline propio (\texttt{.github/workflows/ci.yml}) \\
Arcuri \cite{arcuri2019} & Generación automática de pruebas REST & ACM TOSEM 2019 & Contraste honesto: las 804 pruebas de este proyecto son manuales (JUnit/Mockito), no generadas por \textit{fuzzing} dirigido por búsqueda como EvoMaster --- una vía real de trabajo futuro para ampliar cobertura de ramas sin más esfuerzo manual \\
\bottomrule
\end{tabular}
\caption{Comparison of the thirteen empirical works most directly related to the techniques applied
in this project.}
\label{tab:trabajos-relacionados}
\end{table}
\newpage
